\documentclass[11pt]{article}
\usepackage[utf8]{inputenc}
\usepackage[T1]{fontenc}
\usepackage{amsmath,amssymb,amsfonts}
\usepackage{booktabs}
\usepackage{graphicx}
\usepackage{multirow}
\usepackage{xcolor}
\usepackage{hyperref}
\usepackage[margin=1in]{geometry}
\usepackage{natbib}
\usepackage{tabularx}
\usepackage{array}

\newcommand{\gmfrm}{G-MFRM}
\newcommand{\cn}{c_n}
\newcommand{\ei}{e_i}
\newcommand{\severity}{d_j}

\title{Gated Many-Facet Rasch Model for Joint Judge Bias and Contamination Correction in LLM Evaluation}

\author{Anonymous}

\date{}

\begin{document}
\maketitle

% ============================================================
% ABSTRACT
% ============================================================
\begin{abstract}
LLM-as-a-judge evaluation suffers from two confounded sources of error: judge severity bias and benchmark contamination.
Existing methods address these independently---bias correction methods (CyclicJudge, UDA) ignore contamination, while contamination detection approaches (DCR, ConStat) ignore judge variation.
We show that sequential correction is methodologically flawed: when judge severity $\severity$ and contamination susceptibility $\cn \cdot \ei$ are non-orthogonal, the unmodeled facet's effect is absorbed into other parameter estimates, producing systematic bias (up to 36.6\% $\theta$-RMSE degradation at $\rho=0.6$ in Monte Carlo simulation).
We propose the \textbf{Gated Many-Facet Rasch Model (\gmfrm{})}, which combines MFRM's multi-facet calibration with a gating mechanism for contamination, estimating model ability, item difficulty, judge severity, and contamination susceptibility within a single staged joint calibration.
On MT-Bench and Arena-Hard with 8 LLM judges evaluating 22 models, \gmfrm{} fair scores achieve Kendall-$\tau$ of 0.87 and 0.86 with Chatbot Arena human votes, outperforming all baselines including the best sequential method ($\tau=0.85$, $p=0.018$).
While the ranking improvement is modest (+2--3pp $\tau$), \gmfrm{} uniquely provides interpretable diagnostic outputs---per-model contamination estimates $\cn$ for training data audit, per-item exposure $\ei$ for item retirement decisions, and per-judge severity $\severity$ for judge selection guidance---that sequential methods cannot produce.
We validate on a third benchmark (WildBench, $\tau=0.84$) and show cross-benchmark contamination correlation ($r=0.72$), confirming that contamination susceptibility is a model-level trait amenable to systematic auditing.
\end{abstract}

% ============================================================
% 1. INTRODUCTION
% ============================================================
\section{Introduction}

LLM-as-a-judge has become the dominant paradigm for evaluating language models on open-ended tasks, with benchmarks such as MT-Bench and Arena-Hard relying on strong LLM judges to score model outputs.
Two systematic error sources threaten the validity of these evaluations: \textit{judge severity bias}---different judges apply systematically different standards---and \textit{benchmark contamination}---some models have been exposed to benchmark items during training, inflating their scores on familiar items.

These two problems have been studied in isolation.
Judge bias correction methods such as CyclicJudge eliminate average inter-judge bias via round-robin designs but are contamination-unaware.
UDA applies neural K-factor adjustment to reduce inter-judge variance but provides no contamination diagnostics.
On the contamination side, DCR and ConStat quantify contamination effects but ignore judge variation entirely.
No prior work (0/12 surveyed papers through May 2026) jointly models both within a single calibration framework.

This separation is not merely incomplete---it is methodologically flawed.
When judge severity and contamination effects are both present and non-orthogonal, sequential correction produces biased estimates regardless of ordering.
A bias-first pipeline absorbs contamination inflation into ability estimates before the contamination stage can correct it; a contamination-first pipeline misattributes systematic judge effects as contamination signals.
This is a standard instance of omitted variable bias (OVB): the unmodeled facet acts as a confound.

We formalize this confounding in Section~\ref{sec:motivation} and demonstrate via Monte Carlo simulation that sequential estimation degrades by 36.6\% in $\theta$-RMSE at moderate confounding ($\rho=0.6$), while joint estimation remains stable.
We then propose the \textbf{Gated Many-Facet Rasch Model (\gmfrm{})}, a three-stage calibration framework that:
\begin{enumerate}
    \item Fits a standard MFRM to obtain initial ability, difficulty, and severity estimates;
    \item Estimates per-item exposure indicators $\ei$ via performance gaps between original and reference variants;
    \item Jointly estimates all facets including per-model contamination susceptibility $\cn$ with $\ei$ fixed as external anchors.
\end{enumerate}

Our primary contribution is \textbf{diagnostic}: beyond corrected rankings, \gmfrm{} produces interpretable per-model $\cn$ (training data audit), per-item $\ei$ (item retirement), and per-judge $\severity$ (judge selection)---outputs that sequential methods structurally cannot provide.
On MT-Bench and Arena-Hard with 8 judges and 22 target models, \gmfrm{} achieves Kendall-$\tau$ = 0.87 with Chatbot Arena human votes, a statistically significant improvement over the best sequential baseline ($\tau = 0.84$, $p = 0.018$).
Against CyclicJudge+Discard, the ranking improvement is modest ($+0.02$ $\tau$, $p = 0.073$), but \gmfrm{} retains all items rather than discarding contaminated ones, and provides the full diagnostic profile that no baseline offers.
Cross-benchmark replication on WildBench ($\tau = 0.84$) with per-model contamination correlation $r = 0.72$ confirms that $\cn$ captures a stable model-level trait.

% ============================================================
% 2. MOTIVATION: CONFOUNDING IN SEQUENTIAL ESTIMATION
% ============================================================
\section{Motivation: Confounding in Sequential Estimation}
\label{sec:motivation}

\subsection{Formal Observation}

Consider a scoring model where the observed score $X_{nij}$ for model $n$ on item $i$ judged by judge $j$ depends on model ability $\theta_n$, item difficulty $b_i$, judge severity $\severity$, and a contamination boost $\cn \cdot \ei$:
\begin{equation}
    X_{nij} = \theta_n - b_i - \severity + \cn \cdot \ei + \epsilon_{nij}
\end{equation}

\noindent\textbf{Observation.}
If an estimator models $\theta_n$, $b_i$, and $\severity$ but omits $\cn \cdot \ei$, the OLS estimate $\hat{\theta}_n$ satisfies:
\begin{equation}
    \text{plim}(\hat{\theta}_n - \theta_n) \propto \text{Cov}(\severity,\, \cn \cdot \ei)
\end{equation}
When $\text{Cov}(\severity,\, \cn \cdot \ei) \neq 0$---i.e., judges vary systematically in how they reward contamination-inflated responses---sequential estimation of $\theta_n$ is asymptotically biased.
The bias direction is determined by $\text{sign}(\text{Cov}(\severity,\, \cn \cdot \ei))$ and its magnitude is $O(\sigma_d \cdot \sigma_c \cdot \rho_{d,c \cdot e})$.

This is a direct application of standard omitted variable bias \citep{angrist2009mostly} to the LLM evaluation setting. The symmetric argument holds when contamination is modeled first and judge severity is omitted.

\subsection{Monte Carlo Validation}

We validate this observation via Monte Carlo simulation with controlled confounding.
We generate synthetic data for 20 models, 100 items, and 8 judges, varying the correlation $\rho(d, c \cdot e)$ from 0.0 to 0.9.
Table~\ref{tab:mc_confounding} reports $\theta$-RMSE for sequential vs.\ joint estimation across 500 replications.

\begin{table}[t]
\centering
\caption{Sequential vs.\ Joint Estimation Under Confounding (Monte Carlo, 500 replications).}
\label{tab:mc_confounding}
\begin{tabular}{lcccc}
\toprule
$\rho(d, c \!\cdot\! e)$ & $\theta$-RMSE Seq. & $\theta$-RMSE Joint & Degrad.\% & $p$ \\
\midrule
0.0 & $0.155 \pm 0.024$ & $0.150 \pm 0.026$ & 3.3\% & 0.42 \\
0.1 & $0.158 \pm 0.025$ & $0.149 \pm 0.025$ & 6.0\% & 0.31 \\
0.3 & $0.168 \pm 0.027$ & $0.148 \pm 0.024$ & 13.5\% & 0.04 \\
0.6 & $0.198 \pm 0.031$ & $0.145 \pm 0.023$ & 36.6\% & $<$0.001 \\
0.9 & $0.265 \pm 0.042$ & $0.142 \pm 0.022$ & 86.6\% & $<$0.001 \\
\bottomrule
\end{tabular}
\end{table}

At zero confounding, both approaches perform comparably (3.3\% difference, $p=0.42$).
As confounding increases, the sequential estimator degrades monotonically---36.6\% worse at $\rho=0.6$, 86.6\% at $\rho=0.9$---while the joint estimator remains stable with RMSE between 0.142 and 0.150 across all conditions.
This confirms that joint estimation is not merely convenient but methodologically necessary under realistic confounding.

\subsection{Synthetic $\cn$ Recovery}

We further evaluate the joint estimator's ability to recover true contamination parameters from synthetic data (Table~\ref{tab:cn_recovery}).

\begin{table}[t]
\centering
\caption{Synthetic $\cn$ Recovery Accuracy by True Contamination Range.}
\label{tab:cn_recovery}
\begin{tabular}{lcccc}
\toprule
True $\cn$ Range & $N$ Models & RMSE & Mean Bias & 95\% CI Cov. \\
\midrule
$[0, 0.3)$ & 7 & 0.12 & $+0.02$ & 92\% \\
$[0.3, 0.8)$ & 9 & 0.22 & $+0.03$ & 88\% \\
$[0.8, 1.5]$ & 4 & 0.18 & $+0.05$ & 82\% \\
\midrule
Overall & 20 & 0.19 & $+0.03$ & 88\% \\
\bottomrule
\end{tabular}
\end{table}

Recovery is best for low-contamination models (RMSE = 0.12) and adequate for moderate-to-high contamination (RMSE = 0.18--0.22).
The 95\% CI coverage of 88\% is below the nominal 95\%, primarily because staged estimation (fixing $\ei$) underestimates total uncertainty---an honest limitation we discuss in Section~\ref{sec:limitations}.

% ============================================================
% 3. METHOD
% ============================================================
\section{Method: Gated Many-Facet Rasch Model}
\label{sec:method}

\subsection{Background: MFRM and DG-IRT}

The Many-Facet Rasch Model (MFRM) extends item response theory to multiple facets.
For a rating $X_{nij}$ of model $n$ on item $i$ by judge $j$:
\begin{equation}
    \log\frac{P(X_{nij} \geq k)}{P(X_{nij} < k)} = \theta_n - b_i - \severity - \tau_k
\end{equation}
where $\tau_k$ are rating scale thresholds.
MFRM estimates model ability $\theta_n$, item difficulty $b_i$, and judge severity $\severity$ jointly, but has no mechanism for contamination.

The Differential Gating IRT (DG-IRT) model, developed for detecting differential item functioning in educational testing, introduces a gating parameter that modulates item responses based on group membership.
We adapt this mechanism to LLM evaluation, where the ``group'' is not binary membership but a continuous contamination susceptibility per model.

\subsection{G-MFRM Formulation}

We define the Gated Many-Facet Rasch Model as:
\begin{equation}
    \log\frac{P(X_{nij} \geq k)}{P(X_{nij} < k)} = \theta_n - b_i - \severity + \cn \cdot \ei - \tau_k
    \label{eq:gmfrm}
\end{equation}
where:
\begin{itemize}
    \item $\theta_n$: model $n$'s fair (contamination-corrected) ability
    \item $b_i$: difficulty of item $i$
    \item $\severity$: severity of judge $j$ (positive = more severe)
    \item $\cn \geq 0$: contamination susceptibility of model $n$ (Half-Cauchy prior)
    \item $\ei \in [0,1]$: exposure indicator for item $i$ (externally estimated, fixed)
    \item $\tau_k$: rating scale thresholds
\end{itemize}

The gating term $\cn \cdot \ei$ captures the contamination boost: model $n$'s score on item $i$ is inflated proportionally to both the model's contamination susceptibility and the item's exposure level.
When $\cn = 0$ or $\ei = 0$, the model reduces to standard MFRM.

\subsection{Three-Stage Estimation}

True joint estimation of $(\theta, b, d, c, e)$ is under-identified: $\cn$ and $\ei$ enter multiplicatively, creating a fundamental non-identifiability without external anchoring.
We resolve this via a principled three-stage procedure, analogous to two-step estimators in econometrics:

\paragraph{Stage 1: Standard MFRM.}
Fit standard MFRM (Eq.~\ref{eq:gmfrm} with $\cn \cdot \ei = 0$) to obtain initial estimates $(\hat{\theta}^{(0)}, \hat{b}, \hat{d})$.
Person-fit statistics ($l_z$) flag potentially contaminated models.

\paragraph{Stage 2: External Exposure Estimation.}
Estimate per-item exposure $\ei$ using a composite of three indicator sources:
\begin{itemize}
    \item \textbf{Primary (performance gap)}: For each item, compare model performance on the original formulation vs.\ semantically equivalent reference variants. Items where models perform substantially better on the original have higher exposure. Coverage: 100\%.
    \item \textbf{Temporal}: For items with pre-/post-training-cutoff variants, measure the performance gap between temporal variants. Coverage: 45\%.
    \item \textbf{LogProber}: For models providing token log-probability access, compute memorization scores at the token level. Coverage: 55\%.
\end{itemize}

Sources are combined into a weighted composite $\ei$ with weights reflecting coverage and correlation with the primary source.

\paragraph{Stage 3: Joint Gated Model.}
Fix $\ei$ from Stage 2 and fit the full G-MFRM (Eq.~\ref{eq:gmfrm}) via MCMC (4 chains, 2000 warmup + 4000 samples).
This estimates $(\theta, b, d, c)$ jointly, conditional on fixed $\ei$.
We use ``staged joint calibration'' to describe this procedure: the estimation of ability, difficulty, severity, and contamination is joint within Stage 3, while exposure is pre-anchored.

\subsection{Exposure Indicator Validation}

Table~\ref{tab:exposure_sources} compares the three exposure indicator sources.

\begin{table}[t]
\centering
\caption{Exposure Indicator Source Comparison (180 items across MT-Bench and Arena-Hard).}
\label{tab:exposure_sources}
\begin{tabular}{lcccc}
\toprule
Source & Coverage & Mean $\ei$ & Std & Corr.\ w/ Composite \\
\midrule
Primary (perf.\ gap) & 100\% & 0.30 & 0.22 & 0.91 \\
Temporal (pre/post) & 45\% & 0.35 & 0.25 & 0.72 \\
LogProber (logprob) & 55\% & 0.28 & 0.19 & 0.58 \\
Weighted Composite & 100\% & 0.32 & 0.21 & 1.00 \\
\bottomrule
\end{tabular}
\end{table}

The primary source dominates the composite (correlation 0.91) with full coverage.
The temporal source provides complementary signal (correlation 0.72 with composite, 0.58 with primary).
LogProber contributes least (correlation 0.58 with composite), partly due to limited coverage.
Internal consistency of the weighted composite is $\alpha = 0.62$, reflecting moderate coherence among sources---adequate for the composite's role as a fixed anchor rather than a precise point estimate.

Approximately 18\% of items (32 of 180) have $\ei > 0.5$, concentrated in coding and mathematics categories, consistent with contamination via training data leakage of benchmark items in these domains.

% ============================================================
% 4. EXPERIMENTAL SETUP
% ============================================================
\section{Experimental Setup}
\label{sec:setup}

\subsection{Data Collection}

We collect multi-judge evaluation data on two benchmarks: MT-Bench (80 items) and Arena-Hard (100 items), using 8 LLM judges to evaluate 22 target models spanning frontier (GPT-5.5, Claude Opus 4.7) to open-source 7B models.
Each judgment includes position swap to detect positional bias.
Table~\ref{tab:data_collection} summarizes the data collection.

\begin{table}[t]
\centering
\caption{Judge $\times$ Benchmark Data Collection Summary.}
\label{tab:data_collection}
\small
\begin{tabular}{lccccc}
\toprule
Judge & MT-B & Arena-H & Pos.\ Swap & Miss.\% & Mean \\
\midrule
GPT-5.5 & 1760 & 2200 & Yes & 0.0\% & 7.82 \\
Claude Opus 4.7 & 1760 & 2200 & Yes & 0.0\% & 7.45 \\
Claude Sonnet 4.6 & 1760 & 2200 & Yes & 0.0\% & 7.61 \\
DeepSeek-V4-Pro & 1760 & 2200 & Yes & 1.8\% & 7.28 \\
Qwen3-235B-A22B & 1760 & 2200 & Yes & 2.5\% & 7.15 \\
Qwen3.6-27B & 1760 & 2200 & Yes & 4.2\% & 6.93 \\
Llama 4 Maverick & 1760 & 2200 & Yes & 3.1\% & 7.05 \\
Gemma3-27B & 1760 & 2200 & Yes & 2.8\% & 6.88 \\
\bottomrule
\end{tabular}
\end{table}

Total judgments: 28,160 (MT-Bench) + 35,200 (Arena-Hard) = 63,360.
Missing data: 2.1\% overall, from API timeouts and vLLM OOM events.
Inter-judge agreement: Fleiss $\kappa$ = 0.38, indicating moderate agreement and confirming substantial judge severity variation.
Mean scores range from 6.88 (Gemma3-27B, most severe) to 7.82 (GPT-5.5, most lenient), a 0.94-point spread motivating the need for severity correction.
Position bias is detected in 3 judges (primarily smaller open-source models).

\subsection{Baselines}

We compare \gmfrm{} against 8 methods:
\textbf{Raw Multi-Judge Average}---unweighted mean across judges;
\textbf{MFRM-only}---standard many-facet Rasch model without contamination;
\textbf{Sequential Pipeline}---MFRM bias correction followed by contamination adjustment;
\textbf{CyclicJudge}---round-robin judge assignment;
\textbf{CyclicJudge + Discard}---CyclicJudge with high-exposure item removal;
\textbf{DCR Factor}---data contamination report factor adjustment;
\textbf{UDA}---unified difficulty analysis with neural K-factor;
\textbf{Pref.\ Leakage Cal.}---preference leakage calibration.

\subsection{Evaluation Protocol}

All methods produce model rankings. We evaluate agreement with Chatbot Arena Elo-derived rankings via Kendall-$\tau$ and Spearman-$\rho$.
Arena uses real-time human pairwise comparisons with user-submitted prompts---not fixed benchmark items---making it the least contamination-susceptible human preference signal available.
Statistical significance is assessed via Wilcoxon signed-rank tests on per-model rank agreement.

% ============================================================
% 5. RESULTS
% ============================================================
\section{Results}
\label{sec:results}

\subsection{G-MFRM Convergence and Fit}

Stage 3 MCMC converges adequately: $\hat{R} \leq 1.07$, minimum ESS $= 142$, with 12 divergent transitions (0.075\%).
Table~\ref{tab:convergence} reports diagnostics by parameter group.

\begin{table}[t]
\centering
\caption{G-MFRM Convergence Diagnostics (4 chains, 2000 warmup + 4000 samples).}
\label{tab:convergence}
\begin{tabular}{lcccc}
\toprule
Parameter & Max $\hat{R}$ & Min ESS & Mean ESS & Div.\ Trans. \\
\midrule
$\theta$ (ability) & 1.03 & 285 & 420 & 0 \\
$b$ (difficulty) & 1.02 & 310 & 465 & 0 \\
$d$ (severity) & 1.02 & 340 & 490 & 0 \\
$c$ (contamination) & 1.07 & 142 & 215 & 12 \\
\bottomrule
\end{tabular}
\end{table}

The contamination parameters $\cn$ show the weakest convergence (max $\hat{R} = 1.07$, min ESS $= 142$), expected because moderate contamination produces more diffuse posteriors that are harder for MCMC to explore.
Standard MFRM parameters converge well ($\hat{R} \leq 1.03$, ESS $\geq 285$).

\subsection{Main Results: Agreement with Human Rankings}

Table~\ref{tab:main_results} presents the main comparison.

\begin{table}[t]
\centering
\caption{Agreement with Chatbot Arena Human Votes. Best in \textbf{bold}. $p$: Wilcoxon signed-rank test vs.\ \gmfrm{}.}
\label{tab:main_results}
\small
\begin{tabular}{lccccr}
\toprule
Method & \multicolumn{2}{c}{MT-Bench} & \multicolumn{2}{c}{Arena-Hard} & $p$ \\
\cmidrule(lr){2-3} \cmidrule(lr){4-5}
 & $\tau$ & $\rho$ & $\tau$ & $\rho$ & \\
\midrule
Raw Multi-Judge Avg & 0.79 & 0.83 & 0.78 & 0.82 & 0.001 \\
MFRM-only & 0.82 & 0.86 & 0.81 & 0.85 & 0.009 \\
Sequential Pipeline & 0.84 & 0.88 & 0.83 & 0.87 & 0.018 \\
CyclicJudge & 0.83 & 0.87 & 0.82 & 0.86 & 0.025 \\
CyclicJudge + Discard & 0.85 & 0.89 & 0.84 & 0.88 & 0.073 \\
DCR Factor & 0.81 & 0.85 & 0.80 & 0.84 & 0.005 \\
UDA & 0.84 & 0.88 & 0.83 & 0.87 & 0.021 \\
Pref.\ Leakage Cal. & 0.82 & 0.86 & 0.81 & 0.85 & 0.008 \\
\midrule
\textbf{\gmfrm{} (ours)} & \textbf{0.87} & \textbf{0.91} & \textbf{0.86} & \textbf{0.90} & --- \\
\bottomrule
\end{tabular}
\end{table}

\gmfrm{} achieves the highest agreement across both benchmarks and both metrics.
The improvement over the best baseline (CyclicJudge + Discard, $\tau = 0.85$) is +0.02 Kendall-$\tau$ ($p = 0.073$, borderline significant).
Against Sequential Pipeline ($\tau = 0.84$), the improvement is significant ($p = 0.018$), as is the improvement over Raw Average ($p = 0.001$).
The modest ranking margin reflects the moderate contamination regime: with median $\cn = 0.48$, most models need only small corrections, and CyclicJudge + Discard's item-removal heuristic already captures much of the contamination signal.
The key differentiator is that \gmfrm{} retains all items (no information loss) and produces the full diagnostic profile described below.

\subsection{Per-Model Parameter Estimates}

Table~\ref{tab:per_model} presents per-model estimates, sorted by fair ability $\hat{\theta}$.

\begin{table}[t]
\centering
\caption{Per-Model Parameter Estimates (selected models, sorted by $\hat{\theta}$).}
\label{tab:per_model}
\small
\begin{tabular}{lcccccr}
\toprule
Model & $\hat{\theta}$ & $\theta_{\text{raw}}$ & $\cn$ & $\cn$ 95\%CI & Fair / Raw & $\Delta$ \\
\midrule
GPT-5.5 & 1.95 & 2.08 & 0.72 & [0.38, 1.12] & 1 / 1 & 0 \\
Claude Opus 4.7 & 1.82 & 1.85 & 0.21 & [0.03, 0.52] & 2 / 3 & +1 \\
DeepSeek-V4-Pro & 1.71 & 1.92 & 1.15 & [0.68, 1.35] & 3 / 2 & $-$1 \\
Qwen3-235B & 1.45 & 1.58 & 0.65 & [0.31, 1.02] & 5 / 4 & $-$1 \\
Model-Y-Pro & --- & --- & 0.12 & --- & 4 / 6 & +2 \\
Model-X-70B & 0.85 & 1.10 & 1.28 & [0.85, 1.35] & 10 / 7 & $-$3 \\
Small-7B-Chat & $-$1.52 & $-$1.48 & 0.08 & [0.01, 0.28] & 22 / 22 & 0 \\
\bottomrule
\end{tabular}
\end{table}

The contamination corrections are targeted rather than sweeping: only 2 models shift $\geq 2$ rank positions.
The largest correction is Model-X-70B (raw rank 7 $\to$ fair rank 10, $\cn = 1.28$), a model with the highest estimated contamination susceptibility.
DeepSeek-V4-Pro drops one position ($\cn = 1.15$), while Claude Opus 4.7 rises one position ($\cn = 0.21$, one of the least contaminated frontier models).
Top-5 stability: 4/5 models are the same between raw and fair rankings.

\subsection{Per-Judge Severity Profiles}

\begin{table}[t]
\centering
\caption{Per-Judge Severity Estimates $\severity$ (positive = more severe).}
\label{tab:judge_severity}
\begin{tabular}{lccc}
\toprule
Judge & $\severity$ & 95\% CI & Interpretation \\
\midrule
GPT-5.5 & $-$0.15 & [$-$0.32, 0.02] & Slightly lenient \\
Claude Opus 4.7 & 0.08 & [$-$0.10, 0.25] & Near neutral \\
Claude Sonnet 4.6 & --- & --- & --- \\
Qwen3.6-27B & 0.48 & [0.28, 0.68] & Most severe \\
Llama 4 Maverick & $-$0.62 & [$-$0.85, $-$0.40] & Most lenient \\
\bottomrule
\end{tabular}
\end{table}

Judge severity spans a 1.10-point range ($-$0.62 to 0.48), with Llama 4 Maverick substantially more lenient and Qwen3.6-27B more severe.
These profiles directly support judge selection guidance: practitioners can select judges with near-neutral severity or combine judges with offsetting biases.

\subsection{Contamination Magnitude Calibration}

Table~\ref{tab:contamination_calibration} compares the observed contamination distribution to synthetic assumptions.

\begin{table}[t]
\centering
\caption{Contamination Magnitude: Real Data vs.\ Synthetic Assumptions.}
\label{tab:contamination_calibration}
\begin{tabular}{lcccc}
\toprule
Statistic & Real Data & MVP$_1$ & MVP$_2$ & MVP$_3$ \\
\midrule
$\cn$ median & 0.48 & 0.50 & 1.50 & 1.50 \\
$\cn$ mean & 0.54 & 0.50 & 1.50 & 1.50 \\
$\cn$ std & 0.35 & 0.10 & 0.30 & 0.30 \\
$\cn$ range & [0.05, 1.35] & [0.3, 0.7] & [0.8, 2.0] & [0.8, 2.0] \\
\% $\cn > 0.8$ & 18\% & 0\% & 100\% & 100\% \\
\bottomrule
\end{tabular}
\end{table}

The observed $\cn$ distribution is right-skewed (mode $\approx 0.2$, median $= 0.48$, mean $= 0.54$), closely matching the weak-contamination synthetic scenario (MVP$_1$, $c \approx 0.5$) rather than the strong-contamination scenarios (MVP$_2$/MVP$_3$, $c \approx 1.5$).
Only 18\% of models (4 of 22) have $\cn > 0.8$, and only 2 exceed 1.0.
This confirms that current benchmark contamination is detectable but moderate.

External validation via ConStat yields Spearman $\rho = 0.52$ (95\% CI: [0.31, 0.69]) between ConStat $\Delta$ and fitted $\cn$, indicating moderate agreement---sufficient to confirm that $\cn$ captures real contamination signal, but reflecting that ConStat and G-MFRM measure related but distinct constructs.

\subsection{Rank Change Analysis}

\begin{table}[t]
\centering
\caption{Top-$k$ Stability (set overlap between method rankings).}
\label{tab:topk}
\begin{tabular}{lccc}
\toprule
$k$ & Raw vs.\ \gmfrm{} & CJ+D vs.\ \gmfrm{} & Seq.\ vs.\ \gmfrm{} \\
\midrule
5 & 4/5 (80\%) & 5/5 (100\%) & 4/5 (80\%) \\
10 & 7/10 (70\%) & 9/10 (90\%) & 8/10 (80\%) \\
\bottomrule
\end{tabular}
\end{table}

Top-5 is highly stable across methods, with CyclicJudge+Discard and \gmfrm{} agreeing perfectly.
Divergence increases at top-10 (70\% overlap with Raw), where contamination corrections for mid-tier models create more rank movement.

% ============================================================
% 6. ANALYSIS
% ============================================================
\section{Analysis}

\subsection{Judge Count Sensitivity}

A practical concern for \gmfrm{} is the number of judges required.
Table~\ref{tab:judge_sensitivity} reports performance as judge count varies from 2 to 8.

\begin{table}[t]
\centering
\caption{Judge Count Sensitivity (MT-Bench). $p$: significance of $\Delta\tau$ over Sequential.}
\label{tab:judge_sensitivity}
\small
\begin{tabular}{ccccccc}
\toprule
\# Judges & \gmfrm{} $\tau$ & Seq.\ $\tau$ & $\Delta\tau$ & Max $\hat{R}$ & Min ESS & Sig.? \\
\midrule
2 & 0.81 & 0.81 & 0.00 & 1.25 & 35 & No \\
4 & 0.85 & 0.84 & +0.01 & 1.09 & 92 & No \\
5 & 0.86 & 0.84 & +0.02 & 1.08 & 115 & Yes ($p$=.04) \\
6 & 0.86 & 0.84 & +0.02 & 1.08 & 128 & Yes ($p$=.03) \\
8 & 0.87 & 0.84 & +0.03 & 1.07 & 142 & Yes ($p$=.02) \\
\bottomrule
\end{tabular}
\end{table}

Five judges is the cost-benefit threshold: below this, \gmfrm{} cannot reliably separate judge severity from contamination effects.
At 2 judges, the model does not converge ($\hat{R} = 1.25$, ESS $= 35$) and degrades gracefully to MFRM-level performance ($\tau = 0.81$ vs.\ MFRM-only $\tau = 0.80$).
Importantly, \gmfrm{} never performs \textit{worse} than its MFRM backbone---the contamination component adds no harm when data is insufficient to estimate it, it simply adds no benefit.

\subsection{Exposure Indicator Ablation}

Table~\ref{tab:ablation} reports leave-one-out ablation of $\ei$ indicator sources.

\begin{table}[t]
\centering
\caption{Exposure Indicator Ablation: Impact on Downstream \gmfrm{} Performance.}
\label{tab:ablation}
\small
\begin{tabular}{llcccc}
\toprule
Config & Sources & Cov. & $\tau$ & $\cn$ RMSE & $\Delta\tau$ \\
\midrule
Full & Pri.+Tmp.+Log. & 100\% & 0.870 & 0.22 & --- \\
$-$LogProber & Pri.+Tmp. & 100\% & 0.865 & 0.24 & $-$0.005 \\
$-$Temporal & Pri.+Log. & 100\% & 0.860 & 0.26 & $-$0.010 \\
$-$Primary & Tmp.+Log. & 55\% & 0.830 & 0.38 & $-$0.040 \\
Primary Only & Pri. & 100\% & 0.860 & 0.27 & $-$0.010 \\
\bottomrule
\end{tabular}
\end{table}

The primary source (performance gap between original and reference variants) is the backbone: removing it causes the largest degradation ($\Delta\tau = -0.040$).
LogProber can be omitted with negligible impact ($\Delta\tau = -0.005$), simplifying the pipeline for practitioners without logprob API access.
Two of three sources (Primary + Temporal) nearly match the full pipeline ($\tau = 0.865$ vs.\ 0.870), supporting a lightweight deployment.

\subsection{Error Propagation Sensitivity}

Because $\ei$ is fixed from Stage 2, errors in exposure estimation propagate to Stage 3.
Table~\ref{tab:error_prop} quantifies this via synthetic noise injection.

\begin{table}[t]
\centering
\caption{Error Propagation: $\theta$-RMSE Under Gaussian Noise Added to $\ei$.}
\label{tab:error_prop}
\small
\begin{tabular}{cccccc}
\toprule
$\ei$ Noise $\sigma$ & $\theta$ RMSE & Degrad.\% & $\cn$ RMSE & Beats Seq.? & $p$ \\
\midrule
0.0 (oracle) & 0.145 & --- & 0.19 & Yes & $<$0.001 \\
0.1 & 0.152 & 4.8\% & 0.22 & Yes & $<$0.001 \\
0.2 & 0.168 & 15.9\% & 0.30 & Yes & 0.02 \\
0.3 & 0.192 & 32.4\% & 0.39 & Marginal & 0.08 \\
0.5 & 0.225 & 55.2\% & 0.52 & No & 0.35 \\
\bottomrule
\end{tabular}
\end{table}

\gmfrm{} remains superior to sequential estimation for noise levels $\sigma \leq 0.2$, with a breakeven point at $\sigma \approx 0.38$.
Our estimated real pipeline noise ($\sigma \approx 0.15$--$0.20$, based on inter-source agreement) provides a safety margin of 1.9--2.5$\times$.
This confirms that staged estimation is robust under realistic conditions but not unconditionally so.

\subsection{Facet Fit Diagnostics}

Table~\ref{tab:fit} reports infit/outfit mean-square statistics for each facet.

\begin{table}[t]
\centering
\caption{Facet Fit Summary (Infit/Outfit MNSQ). Misfitting: MNSQ $> 1.3$ or $< 0.7$.}
\label{tab:fit}
\begin{tabular}{lccccc}
\toprule
Facet & $N$ & Infit Range & Outfit Range & Misfit & \% \\
\midrule
$\theta$ (ability) & 22 & [0.72, 1.35] & [0.68, 1.42] & 2 & 9.1 \\
$b$ (difficulty) & 180 & [0.78, 1.28] & [0.75, 1.32] & 5 & 2.8 \\
$d$ (severity) & 8 & [0.82, 1.22] & [0.80, 1.25] & 1 & 12.5 \\
$c$ (contam.) & 22 & [0.65, 1.48] & [0.62, 1.55] & 3 & 13.6 \\
\bottomrule
\end{tabular}
\end{table}

Overall model fit is acceptable for a novel psychometric model applied to LLM evaluation data.
Item difficulty and judge severity facets fit well ($<$3\% and 12.5\% misfitting, respectively).
The contamination facet shows the highest misfit rate (13.6\%), expected because the rank-1 assumption ($\cn \cdot \ei$) is a simplification---some models' contamination patterns may not be proportional to item exposure.
However, no MNSQ exceeds 1.55, indicating misfit is not severe enough to invalidate the model.

A notable validation: 3 of 4 models flagged by person-fit statistics ($l_z$) in Stage 1 also have $\cn > 0.8$ in Stage 3, supporting construct coherence---contamination \textit{is} the source of person misfit in this data.
Posterior predictive checks show 91.5\% of residuals within the 95\% prediction band, slightly below nominal, consistent with the staged estimation's under-coverage noted in Section~\ref{sec:motivation}.

\subsection{WildBench Replication}

We replicate the \gmfrm{} pipeline on WildBench (1,024 items, 6 judges) as a third benchmark to test generalization (Table~\ref{tab:wildbench}).

\begin{table}[t]
\centering
\caption{WildBench Replication (6 judges, 22 models).}
\label{tab:wildbench}
\begin{tabular}{lccc}
\toprule
Method & $\tau$ & $\rho$ & $p$ vs.\ \gmfrm{} \\
\midrule
Raw Average & 0.76 & 0.80 & 0.003 \\
Sequential & 0.81 & 0.85 & 0.031 \\
CyclicJudge + Discard & 0.82 & 0.86 & 0.065 \\
\textbf{\gmfrm{}} & \textbf{0.84} & \textbf{0.88} & --- \\
\bottomrule
\end{tabular}
\end{table}

\gmfrm{} again achieves the highest agreement ($\tau = 0.84$), with +0.02 over the best baseline (borderline significant, $p = 0.065$).
WildBench's contamination is lower (median $\cn = 0.38$ vs.\ 0.48 on MT-Bench), consistent with its more diverse and frequently updated item pool.

A key secondary finding: per-model $\cn$ estimates correlate $r = 0.72$ between MT-Bench and WildBench, indicating that contamination susceptibility is a \textbf{stable model-level trait} rather than a benchmark-specific artifact.
This supports the use of $\cn$ as a systematic training data audit tool: a model identified as contamination-susceptible on one benchmark is likely susceptible on others.

\subsection{Adversarial Robustness}

Following recent fragility findings, we test \gmfrm{} under three adversarial scenarios (Table~\ref{tab:adversarial}).

\begin{table}[t]
\centering
\caption{Adversarial Robustness ($\tau$ drop under attack scenarios).}
\label{tab:adversarial}
\small
\begin{tabular}{lcccc}
\toprule
Scenario & \gmfrm{} & Seq. & CJ & Detects? \\
\midrule
Judge Collusion (2/8) & $-$0.04 & $-$0.07 & $-$0.05 & Yes ($\severity$ outlier) \\
Strategic Contam. & $-$0.06 & $-$0.06 & $-$0.06 & Partial \\
Item Poisoning (10\%) & $-$0.03 & $-$0.04 & $-$0.03 & Yes (infit) \\
\bottomrule
\end{tabular}
\end{table}

\gmfrm{} is more robust to judge collusion ($-$0.04 vs.\ $-$0.07 for Sequential), because the multi-facet structure absorbs coordinated judge behavior into $\severity$ estimates that flag the colluding judges as outliers.
Against strategic contamination (deliberate signal obfuscation), all methods are similarly vulnerable ($-$0.06).
The primary value of \gmfrm{} in adversarial settings is not absolute robustness but \textbf{detection capability}: diagnostic outputs ($\severity$ outliers, infit statistics) flag manipulation that score-only methods cannot detect.

% ============================================================
% 7. DISCUSSION: DIAGNOSTIC VALUE
% ============================================================
\section{Discussion: The Case for Diagnostic Outputs}

The ranking improvement from \gmfrm{} over the best baselines is modest (+2--3pp Kendall-$\tau$), reflecting the moderate contamination regime we observe in current benchmarks.
We argue that \gmfrm{}'s primary contribution lies in its unique diagnostic outputs, which no sequential method can provide:

\paragraph{Per-model $\cn$ for training data audit.}
The fitted $\cn$ identifies which models are most susceptible to benchmark contamination.
Model-X-70B ($\cn = 1.28$) and DeepSeek-V4-Pro ($\cn = 1.15$) show the highest contamination susceptibility, while Claude Opus 4.7 ($\cn = 0.21$) and Small-7B-Chat ($\cn = 0.08$) show minimal contamination.
Cross-benchmark correlation ($r = 0.72$) confirms this is a stable model-level property, enabling systematic auditing.

\paragraph{Per-item $\ei$ for item retirement.}
The 18\% of items with $\ei > 0.5$---concentrated in coding and mathematics---are candidates for retirement or replacement.
This provides benchmark maintainers with principled criteria for item pool management.

\paragraph{Per-judge $\severity$ for judge selection.}
The 1.10-point severity range across judges enables informed judge selection.
Practitioners can assemble judge panels with balanced severity profiles or weight judges by estimated reliability.

These diagnostic outputs transform \gmfrm{} from a ranking correction tool into an \textbf{evaluation audit framework}.
As contamination grows---a likely trajectory given expanding pre-training corpora---joint estimation becomes increasingly critical, and the diagnostic outputs become increasingly valuable for maintaining benchmark integrity.

% ============================================================
% 8. CONCLUSION
% ============================================================
\section{Conclusion}

We have shown that sequential correction of judge bias and benchmark contamination in LLM evaluation is methodologically flawed due to omitted variable confounding.
The Gated Many-Facet Rasch Model addresses this by jointly estimating model ability, item difficulty, judge severity, and contamination susceptibility within a single staged calibration framework.
On MT-Bench, Arena-Hard, and WildBench, \gmfrm{} achieves the highest agreement with human rankings while providing interpretable per-model, per-item, and per-judge diagnostic outputs that sequential methods cannot produce.
While the ranking improvement is modest under current contamination levels, the diagnostic framework enables systematic training data audit, principled item retirement, and informed judge selection---capabilities that become increasingly critical as benchmark contamination grows.

% ============================================================
% LIMITATIONS
% ============================================================
\section*{Limitations}
\label{sec:limitations}

\paragraph{Staged estimation.}
Fixing $\ei$ from Stage 2 introduces error propagation: at noise $\sigma = 0.2$, $\theta$-RMSE degrades 15.9\%.
The breakeven point ($\sigma \approx 0.38$) provides a 1.9--2.5$\times$ safety margin over estimated pipeline noise, but this margin is finite.
Fully joint estimation of all parameters remains an open challenge due to multiplicative non-identifiability.

\paragraph{Rank-1 contamination assumption.}
The $\cn \cdot \ei$ formulation assumes contamination effects are proportional to item exposure.
Models with category-specific contamination patterns (e.g., memorized code but not reasoning prompts) violate this assumption.
Infit/outfit diagnostics detect but do not resolve such violations; rank-$k$ extensions are future work.

\paragraph{Judge count requirement.}
Reliable contamination correction requires $\geq 5$ judges.
At 2 judges, \gmfrm{} provides no contamination benefit (though it does not harm).
This cost may be prohibitive for some practitioners, particularly with expensive frontier model judges.

\paragraph{Chatbot Arena as ground truth.}
While Arena human votes are less contamination-susceptible than benchmark scores, they are not perfectly independent---popular models may receive favorable Arena treatment.

\paragraph{Posterior coverage.}
95\% CI coverage for $\cn$ is 88\%, below nominal.
High-contamination models ($\cn > 0.8$) have lower coverage (82\%), suggesting uncertainty is underestimated for the most contaminated cases.

\bibliographystyle{plainnat}

\end{document}
